% frontmatter/english/actividades.tex -- "The activities": las de
% frontmatter/actividades.tex, en inglés, y "Number names", que solo
% tiene este cuaderno. Tiene que seguir cabiendo en una página
% (tools/check_pages.py comprueba que el día 1 está en la página 5).
\section*{The activities}

Each activity arrives on the day it is first used, and then comes back
from time to time. The instructions on each page, in small grey print,
are read out by a grown-up.

{\small
\begin{multicols}{2}
\raggedright
\subsection*{From autumn}
\begin{itemize}[leftmargin=5mm, itemsep=1pt, topsep=0pt]
\item \textbf{\lblTraza} (on Mondays, when a new number arrives): trace
  the number with a finger over the dots, then with a pencil, and then
  write it alone on the line. The dots are the outline of the real
  number, the same one that is read in the rest of the book.
\item \textbf{\lblColorea}: colour exactly as many as it says, not one
  more -- the rest stay white.
\item \textbf{\lblRodea}: among several groups of things, circle the
  one that has as many as it says. Up to 6, the groups are laid out like
  the dots on a dice, and from 7 to 10 like the ten frame, so they can
  be seen at a glance; the empty group is 0.
\item \textbf{\lblCuenta}: count, pointing to each thing once, and
  circle the number.
\item \textbf{\lblBusca}: circle the number every time it appears,
  among other numbers and shapes.
\item \textbf{\lblUne}: count each group and draw a line from its dot to
  the dot of its number.
\item \textbf{\lblCompleta}: say the numbers on the train out loud, one
  by one, and write the missing ones in the empty carriages.
\item \textbf{\lblDibuja}: draw exactly the things it says, no more and
  no fewer.
\item \textbf{\lblRepasa} (on Fridays): what has been learnt that week,
  to tick, and a drawing. Every ten pages there is a small banner to
  celebrate.
\end{itemize}

\subsection*{From winter}
\begin{itemize}[leftmargin=5mm, itemsep=1pt, topsep=0pt]
\item \textbf{\lblDecena} (11 to 20): a full tray, which is 10, and
  another one next to it; carry on counting from 10.
\item \textbf{\lblCompara}: circle the group with more, or fewer, or the
  two that have the same number; or the biggest number, or the smallest.
\item \textbf{\lblVecinos}: write the number that comes before and the
  one that comes after, in the houses on either side.
\item \textbf{\lblRecta}: say the numbers on the line, from left to
  right, and write the missing ones.
\item \textbf{\lblOrdena}: write the numbers in order, from smallest to
  biggest (or the other way round).
\item \textbf{\lblJunta}: put two groups together and count how many
  there are altogether: adding, before the + and = signs.
\end{itemize}

\subsection*{From spring}
\begin{itemize}[leftmargin=5mm, itemsep=1pt, topsep=0pt]
\item \textbf{\lblSuma} and \textbf{\lblResta}: the same, now with
  their signs; to take away, cross out the ones that go.
\item \textbf{\lblParte}: a line splits the group in two; count what is
  on the other side.
\item \textbf{\lblDiez}: draw the dots that are missing to fill the ten
  frame.
\item \textbf{\lblProblema}: the grown-up reads it out; it can be
  drawn, and the calculation is written down.
\item \textbf{\lblBloques}: each bar is a ten (ten cubes); count the
  bars and the loose cubes.
\item \textbf{\lblNombres}: when 30, 40, 50... arrive, match each number
  to its name -- thir\textbf{teen} or \textbf{thir}ty?
\item \textbf{\lblCompleta}, \textbf{\lblRecta} and \textbf{\lblCuenta},
  also in twos, fives and tens: the wheels of the bikes, the fingers on
  the hands, the tens.
\end{itemize}

\subsection*{From summer}
\begin{itemize}[leftmargin=5mm, itemsep=1pt, topsep=0pt]
\item \textbf{\lblTabla}: each row is a ten; say the numbers row by row,
  and write the missing ones.
\item \textbf{\lblDinero}: count the notes first and then the coins,
  and write how many euros there are (Lucía lives in Spain).
\item \textbf{\lblHora}: the short hand shows the hour; when the long
  hand points straight up, it is “o'clock”, and straight down, “half
  past”.
\item \textbf{\lblMide}: count the cubes under the pencil, from the
  rubber to the tip; with two pencils, circle the longer one.
\end{itemize}
\end{multicols}}
\newpage
